\section{Corpora and Systems}
\label{sec:data}

\paragraph{Primary corpus: Enron e-mail.}
All results in the main text use Enron e-mail~\cite{klimt2004enron} with the questions and gold
answers of the EnronQA benchmark~\cite{ryan2025enronqa}: real e-mail, questions and answers;
tenancy and permissions derived from each message's recipients; a canary sentence injected into
every restricted e-mail.\footnote{\url{https://huggingface.co/datasets/MichaelR207/enron_qa_0922}.}
Ten inboxes, sampled from the 110 that have enough mail, form five tenants of two inboxes each, and
each inbox owner has an assistant. Mail addressed to few people is restricted to its inbox group and
mass mail is internal to the tenant, giving 650 documents of which 500 are restricted (1.7 chunks
each at 300 tokens). E-mails in one thread form a project. Facts are short spans of the gold answers
that occur verbatim in the e-mail and not in the question; 43\% of them also occur in another e-mail
of the corpus, so the judge scores only the corpus-unique ones (\cref{sec:method}).
Following the secret-sharer method~\cite{carlini2019secret}, the canary sentence of each restricted
e-mail is unique, which gives a zero-ambiguity leak signal on real text. The text, questions and
answers are real; the tenancy and ACLs are ours.

\paragraph{Generated corpus: a controlled stub harness.}
A template-generated corpus provides exact ground truth and controlled variation; it backs the test
suite and the pilot study (\cref{app:pilot}) and is not used for the results in the main text. Each seed generates
its own corpus, so seeds are independent replicates rather than re-draws of events from one corpus.
Every corpus has five tenants of 100 documents (40 public, 30 internal, 30 restricted) and eight users
per tenant (six employees, two contractors). Restricted documents come in projects of two, so the twin
world without a target still holds a same-project document, and carry six to twelve filler paragraphs
drawn from sixteen boilerplate memo templates between their fact paragraphs (about 420 words, 2.4 chunks each), so one
document's secrets span several chunks. Numeric fact values are unique across the corpus, every
restricted document carries a canary, and the training tenants (t1--t3) and test tenants (t4--t5) use
disjoint document templates.

\paragraph{Retrieval stack.}
Embeddings come from \texttt{all-MiniLM-L6-v2}~\cite{reimers2019sbert,wang2020minilm}, reranking from
the \texttt{ms-marco-MiniLM-L-6-v2} cross-encoder, and storage from ChromaDB. Retrieval takes the top 5
chunks and passes 3 to the generator, with a minimum cosine relevance of 0.25. The semantic cache uses
a threshold of 0.90 and a TTL of 3600\,s unless varied. Memory writes a summary after 8 turns and keeps
4 turns verbatim.

\paragraph{Generators.}
The model run uses Llama-3.1-8B (instruction-tuned; 4-bit, the default build of Ollama, served
locally) at temperature 0. The same model writes the memory summaries, with a running-summary prompt in the style of
LangChain's summary memory~\cite{doc_langchain_summary}. A deterministic stub that answers
extractively from the best-matching context unit, with an extractive summarizer that keeps sentences
carrying figures, serves as a validation harness and runs the broader defense comparison, the
existence experiment and the post-revocation update experiment; stub and model results are never
pooled.

\paragraph{Clock.}
A simulated clock adds virtual waits (time between probes, padding) to real compute time, so one run
covers up to 48\,h after revocation. The existence experiment uses the stub, whose generation
latency is simulated on the clock; its timing channel therefore reflects the testbed's compute time
and the stub's latency model, not a served model.
